% \begin{table}[!t]
%     \centering
%     \scriptsize
%     \setlength{\tabcolsep}{4pt} % reduces space between columns
%     \begin{tabular}{l cc cc}
%         \toprule
%         \textbf{\# Examples} & \multicolumn{2}{c}{\textbf{Analogy}} & \multicolumn{2}{c}{\textbf{Random}} \\
%         & Coverage & Theorems (avg.) & Coverage  & Theorems (avg.) \\
%         \midrule
%         20 & 88 & 14.21 & 46 & 32.99 \\
%         \midrule
%         50 & 91 & 22.25 & 68 & 50.33 \\
%         \midrule
%         100 & 94 & 30.10 & 78 & 67.33 \\
%          \midrule
%         150 & 96 & 35.68 & 82 & 83.58 \\
%         \bottomrule
%     \end{tabular}
%     \caption{Coverage of 100 target problems using theorems from top-$k$ analogous vs. random examples. Higher $k$ improves coverage but increases the number of theorems in the reduced GDL. Using analogous problems achieves better coverage with fewer theorems. Each target has 3.93 theorems on average; the full GDL contains 196 theorems (not including the variations).
% }
%     \label{tab:theorems_coverage}
% \end{table}



% \begin{table}[!t]
%     \centering
%     \small
%     \setlength{\tabcolsep}{4pt} 
%     \begin{tabular}{l cc cc}
%         \toprule
%         \textbf{$k$} & \multicolumn{2}{c}{\textbf{Analogy}} & \multicolumn{2}{c}{\textbf{Random}} \\
%         & \%Coverage~$\uparrow$ & \#Theorems~$\downarrow$ & \%Coverage~$\uparrow$  & \#Theorems~$\downarrow$  \\
%         \midrule
%         20 & 88 & 11.06 & 62 & 32.92 \\
%         \midrule
%         50 & 93 & 18.57 & 79 & 49.87 \\
%         \midrule
%         100 & 96 & 26.66 & 88 & 66.98 \\
%         \midrule
%         150 & 98 & 31.96 & 92 & 82.82 \\
%         \bottomrule
%     \end{tabular}
%     \caption{Average number of problems that had their entire proof covered by the union of theorems from their top‑$k$ analogies vs.~random problems, over 100 random problems (20 per level, 1–5). Increasing $k$ improves coverage, but results in more theorems. Analogous retrieval consistently achieves higher coverage with fewer theorems.}
%     \label{tab:theorems_coverage}
% \end{table}


\begin{table}[!t]
    \centering
    \small
    \setlength{\tabcolsep}{4pt} 
    \begin{tabular}{l cc cc}
        \toprule
        \textbf{$k$} & \multicolumn{2}{c}{\textbf{Analogy}} & \multicolumn{2}{c}{\textbf{Random}} \\
        & Coverage~$\uparrow$ & Theorems~$\downarrow$ & Coverage~$\uparrow$  & Theorems~$\downarrow$  \\
        \midrule
        20 & 88\% & 11.06 & 62\% & 32.92 \\
        \midrule
        50 & 93\% & 18.57 & 79\% & 49.87 \\
        \midrule
        100 & 96\% & 26.66 & 88\% & 66.98 \\
        \midrule
        150 & 98\% & 31.96 & 92\% & 82.82 \\
        \bottomrule
    \end{tabular}
    \caption{Average number of problems whose entire proof was covered by the union of theorems from their top‑$k$ analogies vs.~random problems, over 100 random target problems (20 per level, 1–5). %Increasing $k$ improves coverage, but results in more theorems. 
    Analogies consistently achieve higher coverage despite fewer theorems.}
    \label{tab:theorems_coverage}
\end{table}



% \dnote{we want to say that we might still be able to prove with those theorems, even if ground truth includes another theorem}

% Moreover, prior work has shown that even models with access to diagrams, such as vision-language models, show only modest gains on similar tasks~\cite{lu2023mathvista}, underscoring the challenge of visual geometric reasoning.


%We then present the top-ranked problems along with their proofs as in-context examples to the LLM, alongside a dictionary of  theorems, expressed in GDL  {(\S\ref{subsec:llm_solver})}.

%Our goal is to prompt the LLM to generate a correct proof for a given target problem.
We use a few-shot prompt (in-context learning) that begins with the target problem, including its textual description, construction, conditions and goal. This is followed by the top analogous problems selected using the regressor from Section~\ref{subsec:similar_problems_retrieval}, along with their full proofs and numeric answers.
%
%
%
Additionally, the LLM is given a theorem dictionary, which defines geometry theorems in terms of formal premises and conclusions (see Figure~\ref{fig:dataset_example}). The LLM's task is to generate a correct proof, using only the theorems from the dictionary. See Appendix~\ref{appendix:LLM_proof_generation} for the prompt.


One challenge we encountered is that the full theorem dictionary contains 196 theorems (234 including  variations), resulting in a large token count and, consequently, 
high costs. This also limits scalability, as adding more theorems would further increase the input size.

To address this problem, we propose a more efficient approach. We have just shown that analogous problems tend to have similar proofs; thus, we test whether we could similarly narrow down the dictionary to include only \emph{theorems used in analogous proofs}. 
This reduction not only reduces costs, but also narrows the search space, helping the model focus on more relevant theorems. 

To evaluate the effectiveness of this approach, we measure the extent to which the narrowed-down dictionary still captures all the theorems needed for the proof of the target problem.
%
We sample 100 problems (20 from each level 1-5) 
and evaluate how many of their target proofs are completely covered using theorems from their top-$k$ similar problems (retrieved with the regressor of  Section~\ref{subsec:similar_problems_retrieval}), and compare to a random set of $k$ problems.

Note that even if a proof of a problem includes a theorem not present in its analogous set, it might still be possible to construct a valid proof using only the covered theorems. %That is, there may be multiple valid ways to prove a given goal. 
Therefore, our coverage metric is a conservative estimate of effectiveness.



See Table~\ref{tab:theorems_coverage} for results.  
Theorems from analogous problems consistently outperform the baseline of theorems from random $k$ problems, providing higher coverage despite lower number of theorems. A larger $k$ increases coverage but also expands the input; we find that $k=100$ offers a good trade-off, achieving coverage of 96\% with only 26.66 theorems on average, that is 13.6\% of the full dictionary. Exploring other values of $k$ is left for future work.





